\documentclass[14pt]{extarticle}
\usepackage{natbib}

\usepackage{xcolor}
\usepackage{amsmath}
\newcommand{\tuple}[1]{ \langle #1 \rangle }
%\usepackage{automata}
\usepackage{times}
\usepackage{ltablex}
\usepackage{tasks}
\usepackage{threeparttable}
\usepackage{booktabs}
\usepackage{xcolor} % For custom colors
\usepackage{tikz} % For styling enumerate numbers
\usepackage{tcolorbox} % For colored box styling
\usepackage{amsmath, amsfonts, amssymb, amsthm} % Math related
\usepackage{natbib}
\usepackage{fontspec}
\usepackage{luatexja}
\usepackage[mathscr]{euscript}

%%%%%% Template
\usepackage{hyperref}
\hypersetup{colorlinks=true,allcolors=blue}

\usepackage{vmargin}
\setpapersize{USletter}
\setmarginsrb{1.0in}{1.0in}{1.0in}{0.6in}{0pt}{0pt}{0pt}{0.4in}

% HOW TO USE THE ABOVE:
%\setmarginsrb{leftmargin}{topmargin}{rightmargin}{bottommargin}{headheight}{headsep}{footheight}{footskip}
%\raggedbottom
% paragraphs indent & skip:
\parindent  0.3cm
\parskip    -0.01cm

\usepackage{tikz}
\usetikzlibrary{backgrounds}

% hyphenation:
% \hyphenpenalty=10000 % no hyphen
% \exhyphenpenalty=10000 % no hyphen
\sloppy

% notes-style paragraph spacing and indentation:
\usepackage{parskip}
\setlength{\parindent}{0cm}

% let derivations break across pages
\allowdisplaybreaks

\newcommand{\orange}[1]{\textcolor{orange}{#1}}
\newcommand{\blue}[1]{\textcolor{blue}{#1}}
\newcommand{\red}[1]{\textcolor{red}{#1}}
\newcommand{\freq}[1]{{\bf \sf F}(#1)}
\newcommand{\datafreq}[2]{{{\bf \sf F}_{#1}(#2)}}

\def\qqquad{\quad\qquad}
\def\qqqquad{\qquad\qquad}

%%%%%%%%%%%%%%%%%%%%%%%%%%%%%%%%%%%%%%%%%%%%%%%%%%%%%%%%%%%%%%%%%%%%%%%%%%%%%%%%
%%%%%%%%%%%%%%%%%%%%%%%%%%%%%%%%%%%%%%%%%%%%%%%%%%%%%%%%%%%%%%%%%%%%%%%%%%%%%%%%

% fill-in-blank question style, found in https://tex.stackexchange.com/a/505089

\usepackage{ifthen}
\usepackage{tocloft}
\usepackage{exercise}
% \usepackage{xcolor}

% Set the Show Answers Boolean
\newboolean{showAns}
\setboolean{showAns}{false}
\newcommand{\showAns}{\setboolean{showAns}{true}}


% ------------------------------------------------------------
% Explanations for each option (only printed when \showAns is enabled)
% Usage: \ExplainChoices{A}{B}{C}{D}
% ------------------------------------------------------------
\newcommand{\ExplainChoices}[4]{%
\ifthenelse{\boolean{showAns}}{%
\vspace{0.3em}\noindent\textcolor{red}{\textbf{Explanations.}}
\begin{itemize}
\item[(A)] #1
\item[(B)] #2
\item[(C)] #3
\item[(D)] #4
\end{itemize}
}{}%
}

% The length of the Answer line
\newlength{\answerlength}
\newcommand{\anslen}[1]{\settowidth{\answerlength}{#1}}

% ans command that indicates space for an answer or shows the answer in red
\newcommand{\ans}[1]{\settowidth{\answerlength}{\hspace{2ex}#1\hspace{2ex}}%
    \ifthenelse{\boolean{showAns}}%
        {\textcolor{red}{\underline{\hspace{2ex}#1\hspace{2ex}}}}%
        {\underline{\hspace{\answerlength}}}}%

\newcommand{\details}[1]{\settowidth{\answerlength}{#1}%
    \ifthenelse{\boolean{showAns}}%
        {\begin{tcolorbox}[colback=blue!5!white,colframe=blue!75!black,title=Solution] #1 \end{tcolorbox}}%
        {}}%

% Formatting how multiple choices Questions are formated.
\settasks{label=(\Alph*), label-width=30pt}


% Some commands for the Exercise Question package
\renewcommand{\QuestionNB}{\Large\protect\textcircled{\small\bfseries\arabic{Question}}\ }
\renewcommand{\ExerciseHeader}{} %no header
\renewcommand{\QuestionBefore}{3ex} %Space above each Q
\setlength{\QuestionIndent}{8pt} % Indent after Q number


% To create the list of answers with tocloft...
\newcommand{\listanswername}{Answers}
\newlistof[Question]{answer}{Answers}{\listanswername}

% Creates a TOC for Answers
\newcounter{prevQ}
\newcommand{\answer}[1]{\refstepcounter{answer}%
\ans{#1}%
\ifnum\theQuestion=\theprevQ%
        \addcontentsline{Answers}{answer}{\protect\numberline{}#1}% don't include the Q number
        \else%
        \addcontentsline{Answers}{answer}{\protect\numberline{\theQuestion}#1}%
        \setcounter{prevQ}{\value{Question}}%
        \fi%
        }%

% \hyphenpenalty=10000 % no hyphen
% \exhyphenpenalty=10000 % no hyphen
\sloppy              % hyphen

\newcommand{\HRule}{\rule{\linewidth}{0.5mm}}
\newcommand{\Hrule}{\rule{\linewidth}{0.3mm}}

%tocloft formatting listofanswers
\renewcommand{\cftAnswerstitlefont}{\bfseries\large}
\renewcommand{\cftanswerdotsep}{\cftnodots}
\cftpagenumbersoff{answer}
\addtolength{\cftanswernumwidth}{10pt}

\makeatletter% since there's an at-sign (@) in the command name
\renewcommand{\@maketitle}{%
  \parindent=0pt% don't indent paragraphs in the title block
  \centering
  {\Large \bfseries\textsc{\@title}} \\
  \vspace{5pt}
  {\large \textit{\@author}} \\
  \HRule \\
  \vspace{1em}
}
\makeatother% resets the meaning of the at-sign (@)


% --- Fonts (robust fallback) ---
\IfFontExistsTF{Crimson Pro Light}{
  \setmainfont{Crimson Pro Light}[
    ItalicFont={* Italic},
    BoldFont={Crimson Pro Medium},
    BoldItalicFont={Crimson Pro Medium Italic}]
  \setsansfont{Crimson Pro Light}[
    ItalicFont={* Italic},
    BoldFont={Crimson Pro Medium},
    BoldItalicFont={Crimson Pro Medium Italic}]
}{
  \setmainfont{HaranoAjiMincho}
  \setsansfont{HaranoAjiGothic}
}
\title{Midterm 2}
\author{Macroeconomics II (Spring 2026)\\ Hui-Jun Chen}


\begin{document}

\maketitle

% Uncomment to display answers in red:
% \showAns

\begin{Exercise}

\section*{Instructions}
\begin{itemize}
\item 40 multiple-choice questions. Exactly one option is correct for each question.
\end{itemize}

\vspace{0.6em}
\Hrule
\vspace{0.8em}

% ============================================================
\section*{Problem 1 (Questions 1--20): Building blocks and key derivations}
% ============================================================

\Question \textbf{Three elements.} The IS--MP--PC model combines:
\answer{C}
\begin{tasks}(1)
\task AD curve, AS curve, and a money demand curve
\task IS curve, LM curve, and a Phillips curve
\task Phillips curve, IS curve, and a monetary policy rule
\task Phillips curve, money supply rule, and a government budget constraint
\end{tasks}
\ExplainChoices{Wrong: money demand/L(M/P) belongs to IS–LM; IS–MP replaces money-market closure with a policy rule.}{Wrong: IS–LM includes LM, not MP; IS–MP–PC replaces LM with MP.}{Correct: the model is demand (IS) + policy rule (MP) + inflation block (PC).}{Wrong: a government budget constraint is not one of the three blocks in IS–MP–PC.}



\Question \textbf{Phillips curve (form).} The baseline Phillips curve is:
\answer{A}
\begin{tasks}(1)
\task $\pi_t=\pi_t^{e}+\gamma (y_t-y_t^\ast)+\varepsilon_t^\pi$
\task $\pi_t=\pi_t^{e}-\gamma (y_t-y_t^\ast)+\varepsilon_t^\pi$
\task $\pi_t=\pi^\ast+\gamma (y_t-y_t^\ast)$
\task $\pi_t=\pi_{t-1}+\gamma(y_t-y_t^\ast)$
\end{tasks}
\ExplainChoices{Correct: inflation equals expected inflation + slope×output gap + a cost-push shock.}{Wrong sign: higher output gap should raise inflation, not lower it.}{Wrong: would ignore expectations and shocks, forcing inflation toward target mechanically.}{Wrong: accelerationist/adaptive form; not the baseline expectations-augmented PC.}



\Question \textbf{Interpretation of $\pi_t^e$.} In the Phillips curve, $\pi_t^e$ represents:
\answer{B}
\begin{tasks}(1)
\task realized inflation last period
\task expected inflation built into wage/price setting at time $t$
\task the inflation target set by the central bank
\task the nominal interest rate
\end{tasks}
\ExplainChoices{Wrong: that is realized past inflation only under adaptive expectations.}{Correct: expected inflation built into wage/price setting at time $t$.}{Wrong: $\pi^*$ is the target, not private expectations.}{Wrong: $i_t$ is the nominal rate.}



\Question \textbf{Transmission mechanism (review emphasis).} A typical chain from inflation to real activity under an inflation-targeting rule is:
\answer{D}
\begin{tasks}(1)
\task $\pi\uparrow \Rightarrow u\downarrow \Rightarrow x\downarrow$
\task $\pi\uparrow \Rightarrow r\downarrow \Rightarrow x\downarrow$
\task $\pi\uparrow \Rightarrow y^\ast\uparrow \Rightarrow x\uparrow$
\task $\pi\uparrow \Rightarrow$ tighter policy $\Rightarrow u\uparrow \Rightarrow x\downarrow$
\end{tasks}
\ExplainChoices{Wrong: tightening raises unemployment; it does not fall in the stabilizing story.}{Wrong: with inflation targeting, inflation above target tends to raise the real rate, not lower it.}{Wrong: inflation does not mechanically raise potential output.}{Correct: $\pi\uparrow$ → tighter policy → unemployment up → output gap down.}



\Question \textbf{Okun connection.} If $u_t-u_t^n\approx -\lambda (y_t-y_t^\ast)$ with $\lambda>0$, then $y_t>y_t^\ast$ implies:
\answer{A}
\begin{tasks}(1)
\task $u_t<u_t^n$
\task $u_t>u_t^n$
\task $u_t=u_t^n$
\task no implication for unemployment
\end{tasks}
\ExplainChoices{Correct: positive output gap implies unemployment below natural ($u<u^n$).}{Wrong: that would correspond to a negative output gap.}{Wrong: equality holds only when the gap is zero.}{Wrong: Okun implies a sign relationship when $\lambda>0$.}



\Question \textbf{IS curve (form).} The IS curve used in the notes is:
\answer{C}
\begin{tasks}(1)
\task $y_t=y_t^\ast+\alpha(i_t-\pi_t-r^\ast)+\varepsilon_t^y$
\task $y_t=y_t^\ast-\alpha(i_t+\pi_t-r^\ast)+\varepsilon_t^y$
\task $y_t=y_t^\ast-\alpha(i_t-\pi_t-r^\ast)+\varepsilon_t^y$
\task $y_t=r^\ast-\alpha(i_t-\pi_t-y_t^\ast)+\varepsilon_t^y$
\end{tasks}
\ExplainChoices{Wrong sign: higher real-rate gap should reduce output, not raise it.}{Wrong: uses $i+\pi$ rather than the real rate $i-\pi$.}{Correct: $y=y^* -\alpha(i-\pi-r^*)+\varepsilon^y$.}{Wrong: incorrect rearrangement; not an IS curve.}



\Question \textbf{What does $\alpha$ measure?} In the IS curve, a larger $\alpha$ means:
\answer{B}
\begin{tasks}(1)
\task output responds less to the real interest rate
\task output responds more to the real interest rate
\task inflation responds more to output
\task the central bank responds more to inflation
\end{tasks}
\ExplainChoices{Wrong: larger $\alpha$ means more sensitivity, not less.}{Correct: output responds more to the real interest-rate gap.}{Wrong: inflation sensitivity is $\gamma$ in the Phillips curve.}{Wrong: policy aggressiveness is $\beta_\pi$.}



\Question \textbf{Monetary policy rule (baseline).} The baseline MP rule is:
\answer{D}
\begin{tasks}(1)
\task $i_t=r^\ast+\pi^\ast-\beta_\pi(\pi_t-\pi^\ast)$
\task $i_t=\pi_t+\beta_\pi(\pi_t-\pi^\ast)$
\task $i_t=r^\ast+\pi_t+\beta_\pi(\pi_t-\pi^\ast)$
\task $i_t=r^\ast+\pi^\ast+\beta_\pi(\pi_t-\pi^\ast)$
\end{tasks}
\ExplainChoices{Wrong sign: would cut rates when inflation rises above target.}{Wrong: missing intercept; not the baseline rule.}{Wrong: double-counts inflation level; baseline uses deviation $(\pi-\pi^*)$.}{Correct: $i=r^*+\pi^*+\beta_\pi(\pi-\pi^*)$.}



\Question \textbf{What does $\beta_\pi$ do?} A higher $\beta_\pi$ implies:
\answer{A}
\begin{tasks}(1)
\task a stronger increase in $i_t$ when inflation rises above target
\task a weaker increase in $i_t$ when inflation rises above target
\task a steeper Phillips curve
\task a higher natural output $y_t^\ast$
\end{tasks}
\ExplainChoices{Correct: stronger nominal-rate response to inflation deviations.}{Wrong: weaker response corresponds to smaller $\beta_\pi$.}{Wrong: Phillips slope is $\gamma$.}{Wrong: $y^*$ is supply-side, not set by $\beta_\pi$.}



\Question \textbf{Special case $\beta_\pi=1$.} Under the baseline MP rule, if $\beta_\pi=1$ then:
\answer{C}
\begin{tasks}(1)
\task the nominal rate does not react to inflation at all
\task the real rate rises when inflation rises
\task the real rate is pinned at $r^\ast$ regardless of $\pi_t-\pi^\ast$
\task the IS--MP curve is downward sloping
\end{tasks}
\ExplainChoices{Wrong: $\beta_\pi=1$ still reacts one-for-one; $\beta_\pi=0$ is no reaction.}{Wrong: the real rate does not rise with inflation deviations when $\beta_\pi=1$.}{Correct: $i=r^*+\pi$ so $i-\pi=r^*$ regardless of $(\pi-\pi^*)$.}{Wrong: downward IS–MP slope requires $\beta_\pi>1$.}



\Question \textbf{Derivation step (review emphasis).} The IS--MP curve is obtained by:
\answer{B}
\begin{tasks}(1)
\task substituting the Phillips curve into the MP rule
\task substituting the MP rule into the IS curve
\task substituting money demand into the IS curve
\task substituting the IS curve into the Phillips curve
\end{tasks}
\ExplainChoices{Wrong: does not yield IS–MP.}{Correct: substitute the MP rule into the IS curve to eliminate $i$.}{Wrong: money demand is not part of IS–MP–PC.}{Wrong: this would not isolate output as a function of inflation via policy.}



\Question \textbf{IS--MP curve (result).} The derived IS--MP curve is:
\answer{A}
\begin{tasks}(1)
\task $y_t=y_t^\ast-\alpha(\beta_\pi-1)(\pi_t-\pi^\ast)+\varepsilon_t^y$
\task $y_t=y_t^\ast+\alpha(\beta_\pi-1)(\pi_t-\pi^\ast)+\varepsilon_t^y$
\task $y_t=y_t^\ast-\alpha\beta_\pi(\pi_t-\pi^\ast)+\varepsilon_t^y$
\task $y_t=y_t^\ast-\alpha(\pi_t-\pi^\ast)+\varepsilon_t^\pi$
\end{tasks}
\ExplainChoices{Correct: $y=y^* -\alpha(\beta_\pi-1)(\pi-\pi^*)+\varepsilon^y$.}{Wrong: sign would imply inflation above target raises output in the Taylor-principle case.}{Wrong: should be $(\beta_\pi-1)$ due to real-rate gap algebra.}{Wrong: $\varepsilon^\pi$ belongs to PC, not IS–MP.}



\Question \textbf{Slope sign.} In a graph with output on the horizontal axis and inflation on the vertical axis, the IS--MP curve slopes down when:
\answer{D}
\begin{tasks}(1)
\task $\beta_\pi<1$
\task $\beta_\pi=0$
\task $\beta_\pi<0$
\task $\beta_\pi>1$
\end{tasks}
\ExplainChoices{Wrong: $\beta_\pi<1$ gives upward slope (destabilizing).}{Wrong: $\beta_\pi=0$ is a peg; slope depends on $(\beta_\pi-1)$.}{Wrong: not the relevant case; sign logic is about above/below 1.}{Correct: $\beta_\pi>1$ → real rate rises with inflation → output falls → downward slope.}



\Question \textbf{Why $\beta_\pi>1$ matters (core intuition).} The Taylor principle is central because when inflation rises it should:
\answer{B}
\begin{tasks}(1)
\task lower the real interest rate, boosting demand
\task raise the real interest rate, reducing demand
\task raise natural output $y^\ast$
\task make the Phillips curve flat
\end{tasks}
\ExplainChoices{Wrong: that amplifies inflation by lowering real rates.}{Correct: stabilizing policy requires raising the real rate when inflation rises.}{Wrong: potential output is not pinned by this policy parameter.}{Wrong: policy does not mechanically change the Phillips slope.}



\Question \textbf{PC shift: higher expected inflation.} Holding $y_t-y_t^\ast$ fixed, an increase in $\pi_t^e$ causes:
\answer{A}
\begin{tasks}(1)
\task $\pi_t$ to rise one-for-one (PC shifts up)
\task $\pi_t$ to fall one-for-one (PC shifts down)
\task no change in $\pi_t$
\task the slope $\gamma$ to fall
\end{tasks}
\ExplainChoices{Correct: higher $\pi^e$ raises $\pi$ one-for-one holding the gap fixed.}{Wrong: it does not shift down.}{Wrong: it changes the intercept, so inflation changes.}{Wrong: slope $\gamma$ is unchanged; intercept shifts.}



\Question \textbf{After $\pi^e$ rises (qualitative result).} If $\pi_t^e$ rises above the target $\pi^\ast$, then in the one-period IS--MP--PC intersection (with $\beta_\pi>1$) we typically get:
\answer{C}
\begin{tasks}(1)
\task higher output and higher inflation, and inflation exceeds expectations
\task lower output and lower inflation
\task lower output and higher inflation, but inflation rises less than expectations
\task higher output and lower inflation
\end{tasks}
\ExplainChoices{Wrong: inflation is not typically above expectations in this linear form.}{Wrong: higher $\pi^e$ shifts PC up, so inflation won’t fall.}{Correct: inflation rises and output falls; with $\beta_\pi>1$, $\Delta\pi<\Delta\pi^e$ (partial pass-through).}{Wrong: higher expectations do not lower inflation here.}



\Question \textbf{Soft vs tough central bank (definition).} In these notes, a ``tough'' central bank corresponds to:
\answer{B}
\begin{tasks}(1)
\task smaller $\beta_\pi$
\task larger $\beta_\pi$
\task smaller $\gamma$
\task larger $\varepsilon_t^\pi$
\end{tasks}
\ExplainChoices{Wrong: smaller $\beta_\pi$ is softer.}{Correct: larger $\beta_\pi$ is tougher.}{Wrong: $\gamma$ is Phillips slope.}{Wrong: $\varepsilon^\pi$ is a shock, not policy.}



\Question \textbf{Tough CB (one-period tradeoff).} When $\beta_\pi$ is higher (tougher CB), a cost-push increase in inflation tends to produce:
\answer{D}
\begin{tasks}(1)
\task a larger rise in inflation and a smaller fall in output
\task a larger rise in inflation and a larger fall in output
\task a smaller rise in inflation and a smaller fall in output
\task a smaller rise in inflation and a larger fall in output
\end{tasks}
\ExplainChoices{Wrong: tougher policy tends to reduce the inflation increase, not raise it.}{Wrong: inflation rise is smaller, not larger, under tougher policy.}{Wrong: output loss is typically larger under tougher policy.}{Correct: smaller inflation rise but larger output fall.}



\Question \textbf{Numerical slope check.} Suppose $\alpha=2$ and $\beta_\pi=2$. The coefficient on $(\pi_t-\pi^\ast)$ in the IS--MP curve is:
\answer{A}
\begin{tasks}(1)
\task $-2$
\task $-4$
\task $+2$
\task $0$
\end{tasks}
\ExplainChoices{Correct: $-\alpha(\beta_\pi-1)=-2(1)=-2$.}{Wrong: would require $(\beta_\pi-1)=2$.}{Wrong: sign should be negative when $\beta_\pi>1$.}{Wrong: zero only when $\beta_\pi=1$.}



\Question \textbf{Interpretation.} In the IS--MP curve, a positive $\varepsilon_t^y$ shifts the curve:
\answer{B}
\begin{tasks}(1)
\task left (lower $y$ for each $\pi$)
\task right (higher $y$ for each $\pi$)
\task up (higher $\pi$ for each $y$)
\task down (lower $\pi$ for each $y$)
\end{tasks}
\ExplainChoices{Correct: $(\pi_t-\pi_{t-1})$ is increasing in the output gap (accelerationist).}{Wrong: the level depends on past inflation and shocks, not only gap.}{Wrong: target alone does not govern the change in inflation.}{Wrong: money growth is not in this reduced form.}



\vspace{0.6em}
\Hrule
\vspace{0.8em}

% ============================================================
\section*{Problem 2 (Questions 21--40): Solving the model and the Taylor principle}
% ============================================================

\Question \textbf{Solving the two-equation system.} After deriving IS--MP, equilibrium inflation is solved by combining:
\answer{C}
\begin{tasks}(1)
\task IS and MP only
\task MP and money demand
\task IS--MP and Phillips curve
\task Phillips curve and money demand
\end{tasks}
\ExplainChoices{Wrong: IS–MP alone does not pin down inflation.}{Wrong: money demand is not in IS–MP–PC.}{Correct: combine IS–MP with the Phillips curve.}{Wrong: money demand is not used here.}



\Question \textbf{Solve for inflation (structure).} Solving yields a formula of the form:
\answer{A}
\begin{tasks}(1)
\task $\pi_t=\theta\pi_t^e+(1-\theta)\pi^\ast+\theta(\text{shocks})$
\task $\pi_t=\pi_t^e+\pi^\ast+\theta(\text{shocks})$
\task $\pi_t=\pi^\ast$ always
\task $\pi_t=\theta\pi^\ast+(1-\theta)\pi_t^e$
\end{tasks}
\ExplainChoices{Correct: convex combination of expectations and target plus shocks.}{Wrong: not additive like that; weights matter.}{Wrong: would eliminate fluctuations.}{Wrong: missing the shock term and correct weights.}



\Question \textbf{Definition of $\theta$.} In the Taylor principle notes:
\answer{B}
\begin{tasks}(1)
\task $\theta = 1+\alpha\gamma(\beta_\pi-1)$
\task $\theta = \dfrac{1}{1+\alpha\gamma(\beta_\pi-1)}$
\task $\theta = \dfrac{1}{1-\alpha\gamma(\beta_\pi-1)}$
\task $\theta = \alpha\gamma(\beta_\pi-1)$
\end{tasks}
\ExplainChoices{Wrong: that is the denominator, not $\theta$.}{Correct: $\theta=1/(1+\alpha\gamma(\beta_\pi-1))$.}{Wrong: sign in denominator is wrong.}{Wrong: that’s only part of the expression.}



\Question \textbf{Case $\beta_\pi>1$.} If $\beta_\pi>1$ and $\alpha,\gamma>0$, then:
\answer{C}
\begin{tasks}(1)
\task $\theta>1$
\task $\theta<0$
\task $0<\theta<1$
\task $\theta=1$
\end{tasks}
\ExplainChoices{Wrong: $\theta>1$ occurs when denominator is between 0 and 1.}{Wrong: $\theta$ is positive here.}{Correct: denominator exceeds 1 so $0<\theta<1$.}{Wrong: $\theta=1$ only if $(\beta_\pi-1)=0$ or $\alpha\gamma=0$.}



\Question \textbf{Parameter interpretation (review emphasis).} In $\theta=\dfrac{1}{1+\alpha\gamma(\beta_\pi-1)}$, $\theta$ is smaller when:
\answer{D}
\begin{tasks}(1)
\task $\alpha$ is smaller
\task $\gamma$ is smaller
\task $\beta_\pi$ is closer to 1
\task $\alpha$, $\gamma$, or $\beta_\pi$ is larger (with $\beta_\pi>1$)
\end{tasks}
\ExplainChoices{Wrong: smaller $\alpha$ increases $\theta$.}{Wrong: smaller $\gamma$ increases $\theta$.}{Wrong: $\beta_\pi$ closer to 1 from above increases $\theta$.}{Correct: larger $\alpha$, $\gamma$, or $\beta_\pi$ (with $\beta_\pi>1$) lowers $\theta$.}



\Question \textbf{Numerical $\theta$.} Let $\alpha=1$, $\gamma=0.5$, $\beta_\pi=3$. Then $\theta$ equals:
\answer{B}
\begin{tasks}(1)
\task $1/(1+0.5\cdot 3)=0.4$
\task $1/(1+0.5\cdot 2)=0.5$
\task $1/(1+0.5\cdot 1)=0.667$
\task $1/(1-0.5\cdot 2)=\infty$
\end{tasks}
\ExplainChoices{Wrong: uses $\beta_\pi$ instead of $(\beta_\pi-1)$.}{Correct: $\theta=1/(1+0.5\cdot2)=0.5$.}{Wrong: would correspond to $(\beta_\pi-1)=1$.}{Wrong: denominator is not zero here.}



\Question \textbf{Effect of tougher policy on inflation.} With $\beta_\pi>1$, raising $\beta_\pi$ tends to:
\answer{A}
\begin{tasks}(1)
\task reduce $\theta$ and make inflation closer to the target
\task increase $\theta$ and make inflation closer to expectations
\task reduce $\gamma$
\task shift the Phillips curve down permanently
\end{tasks}
\ExplainChoices{Correct: higher $\beta_\pi$ lowers $\theta$, anchoring inflation closer to target.}{Wrong: in stable region, $\theta$ falls, not rises.}{Wrong: $\gamma$ is not altered by policy in this reduced form.}{Wrong: policy does not permanently shift the Phillips curve.}



\Question \textbf{Adaptive expectations.} Under adaptive expectations:
\answer{B}
\begin{tasks}(1)
\task $\pi_t^e=\pi^\ast$
\task $\pi_t^e=\pi_{t-1}$
\task $\pi_t^e=\mathbb{E}_t\pi_{t+1}$
\task $\pi_t^e=0$ always
\end{tasks}
\ExplainChoices{Wrong: that is perfectly anchored at target.}{Correct: $\pi_t^e=\pi_{t-1}$.}{Wrong: forward-looking uses $\mathbb{E}_t\pi_{t+1}$.}{Wrong: not the adaptive rule.}



\Question \textbf{Accelerationist form.} With adaptive expectations, the Phillips curve becomes:
\answer{C}
\begin{tasks}(1)
\task $\pi_t=\pi^\ast+\gamma(y_t-y_t^\ast)+\varepsilon_t^\pi$
\task $\pi_t=\pi_t^e-\gamma(y_t-y_t^\ast)+\varepsilon_t^\pi$
\task $\pi_t=\pi_{t-1}+\gamma(y_t-y_t^\ast)+\varepsilon_t^\pi$
\task $\pi_t=\pi_{t-1}-\gamma(y_t-y_t^\ast)+\varepsilon_t^\pi$
\end{tasks}
\ExplainChoices{Wrong: missing lagged inflation term under adaptive expectations.}{Wrong: sign error on gap term.}{Correct: $\pi_t=\pi_{t-1}+\gamma(y-y^*)+\varepsilon^\pi$.}{Wrong: wrong sign on the gap term.}



\Question \textbf{Interpretation.} Under adaptive expectations, there is a positive relationship between:
\answer{A}
\begin{tasks}(1)
\task the change in inflation $(\pi_t-\pi_{t-1})$ and the output gap
\task the level of inflation and the output gap only
\task the output gap and the inflation target only
\task money growth and the output gap only
\end{tasks}
\ExplainChoices{Correct: $(\pi_t-\pi_{t-1})$ is increasing in the output gap (accelerationist).}{Wrong: the level depends on past inflation and shocks, not only gap.}{Wrong: target alone does not govern the change in inflation.}{Wrong: money growth is not in this reduced form.}



\Question \textbf{Stability under Taylor principle.} If $0<\theta<1$, then inflation dynamics $\pi_t=\theta\pi_{t-1}+\cdots$ are:
\answer{D}
\begin{tasks}(1)
\task explosive
\task oscillatory with increasing amplitude
\task a random walk
\task stable AR(1) mean reversion
\end{tasks}
\ExplainChoices{Wrong: explosive would require coefficient >1 in absolute value.}{Wrong: oscillation would require negative coefficient with |.|>1, etc.}{Wrong: random walk corresponds to coefficient 1.}{Correct: $0<\theta<1$ implies stable AR(1) mean reversion.}



\Question \textbf{When $\beta_\pi<1$.} If $\beta_\pi<1$ (so policy violates Taylor principle), then inflation shocks tend to:
\answer{B}
\begin{tasks}(1)
\task raise the real rate and reduce output, stabilizing inflation
\task lower the real rate and raise output, amplifying inflation
\task not affect the real rate
\task stabilize inflation faster
\end{tasks}
\ExplainChoices{Wrong: that is the stabilizing Taylor-principle case.}{Correct: real rate falls when inflation rises → demand increases → amplifies inflation.}{Wrong: real rate changes; it falls in this case.}{Wrong: does not stabilize faster; it destabilizes.}



\Question \textbf{Boundary emphasized in TaylorRule notes.} The interval producing $\theta>1$ is:
\answer{A}
\begin{tasks}(1)
\task $\left(1-\frac{1}{\alpha\gamma}\right)<\beta_\pi<1$
\task $\beta_\pi>1$
\task $\beta_\pi<\left(1-\frac{1}{\alpha\gamma}\right)$
\task $\beta_\pi=1$
\end{tasks}
\ExplainChoices{Correct: $(1-1/(\alpha\gamma),1)$ produces $\theta>1$ (denominator between 0 and 1).}{Wrong: $\beta_\pi>1$ gives $0<\theta<1$.}{Wrong: below $1-1/(\alpha\gamma)$ the denominator becomes negative (different pathology).}{Wrong: $\beta_\pi=1$ gives $\theta=1$.}



\Question \textbf{Numerical boundary.} If $\alpha=1$ and $\gamma=0.5$, then $1-\frac{1}{\alpha\gamma}$ equals:
\answer{C}
\begin{tasks}(1)
\task $0$
\task $0.5$
\task $-1$
\task $2$
\end{tasks}
\ExplainChoices{Wrong: not correct for $\alpha\gamma=0.5$.}{Wrong: not correct; $1-2=-1$.}{Correct: $1-1/0.5=-1$.}{Wrong: sign error.}



\Question \textbf{Graph logic (review emphasis).} When $\beta_\pi<1$, the IS--MP curve in $(y,\pi)$ space:
\answer{B}
\begin{tasks}(1)
\task slopes down
\task slopes up
\task is horizontal
\task is vertical
\end{tasks}
\ExplainChoices{Wrong: downward slope is $\beta_\pi>1$.}{Correct: $\beta_\pi<1$ implies upward-sloping IS–MP in $(y,\pi)$.}{Wrong: not implied.}{Wrong: not implied.}



\Question \textbf{Condition ``upward and steeper than PC''.} The TaylorRule notes state that IS--MP slopes up and is steeper than PC when:
\answer{D}
\begin{tasks}(1)
\task $\beta_\pi>1$
\task $\beta_\pi<1-\frac{1}{\alpha\gamma}$
\task $\beta_\pi=1$
\task $1-\frac{1}{\alpha\gamma}<\beta_\pi<1$
\end{tasks}
\ExplainChoices{Wrong: $\beta_\pi>1$ gives downward IS–MP.}{Wrong: below $1-1/(\alpha\gamma)$ the denominator is negative; not the “steeper than PC” case emphasized.}{Wrong: $\beta_\pi=1$ makes IS–MP vertical in $(y,\pi)$ (output pinned).}{Correct: $1-1/(\alpha\gamma)<\beta_\pi<1$ yields upward IS–MP steeper than PC → explosive dynamics (adaptive expectations).}



\Question \textbf{Explosive dynamics intuition.} In the explosive region $1-\frac{1}{\alpha\gamma}<\beta_\pi<1$, an increase in expected inflation $\pi_t^e$ tends to:
\answer{A}
\begin{tasks}(1)
\task start a sequence where inflation and output move further away from target over time
\task reduce inflation and raise output back to natural values
\task have no effect because expectations do not matter
\task immediately pin inflation at the target
\end{tasks}
\ExplainChoices{Correct: higher expected inflation shifts PC up; with destabilizing IS–MP feedback the economy moves further away over time.}{Wrong: describes stable reversion.}{Wrong: expectations matter via the PC intercept.}{Wrong: inflation is not pinned at target here.}



\Question \textbf{A puzzling result (TaylorRule slide 11).} When $\beta_\pi<1$, increasing the inflation target $\pi^\ast$ can lead to:
\answer{C}
\begin{tasks}(1)
\task higher inflation and higher output always
\task no change in nominal interest rates
\task higher interest rates and lower inflation (counterintuitive outcome)
\task lower interest rates and higher inflation
\end{tasks}
\ExplainChoices{Wrong: not “always”; the slide highlights a paradoxical comparative static under $\beta_\pi<1$.}{Wrong: nominal rates change with the rule/target.}{Correct: the notes highlight a case where raising $\pi^*$ can imply higher rates and lower inflation (counterintuitive).}{Wrong: this is the usual intuition, but the slide emphasizes a reversal can occur.}



\Question \textbf{Quick computation (IS--MP coefficient).} If $\alpha=1.5$ and $\beta_\pi=0.8$, then $-\alpha(\beta_\pi-1)$ equals:
\answer{B}
\begin{tasks}(1)
\task $-0.3$
\task $+0.3$
\task $-1.2$
\task $+1.2$
\end{tasks}
\ExplainChoices{Wrong: sign error.}{Correct: $-1.5(0.8-1)= -1.5(-0.2)=+0.3$.}{Wrong: magnitude incorrect.}{Wrong: magnitude incorrect.}



\Question \textbf{Policy interpretation.} If $-\alpha(\beta_\pi-1)>0$, then higher inflation is associated with:
\answer{A}
\begin{tasks}(1)
\task higher output along IS--MP (destabilizing feedback)
\task lower output along IS--MP (stabilizing feedback)
\task no change in output
\task lower inflation
\end{tasks}
\ExplainChoices{Correct: positive slope means higher inflation is associated with higher output along IS–MP (destabilizing feedback).}{Wrong: stabilizing requires negative slope (higher inflation → lower output).}{Wrong: output changes with inflation when slope ≠ 0.}{Wrong: this is about output response, not a direct statement that inflation falls.}



\end{Exercise}

\end{document}
